\chapter{Locally conformally K\"ahler geometry}
\label{chap:lck}

Locally conformally K\"ahler geometry gives a useful testing ground for the theory developed here. Its local K\"ahler metrics make some geometric constructions available, while its global Lee form records an obstruction to patching those metrics without changing scale. The Hopf example of \cref{chap:hopf} belongs to this class. We explain both its additional structure and the distinction between ordinary and twisted potentials. The final section reconstructs a bimeromorphic invariant from Gauduchon tests and gives an explicit example showing the limits of that invariant. We follow the conventions of this book; the operator $\dc$ in \cite{OV24LCK} differs by a constant factor.

\section{The Lee form and the K\"ahler covering}

Let $X$ be a connected compact complex manifold of dimension $n\geq2$.

\begin{definition}\label{def:lck}
A Hermitian form $\omega$ is \emph{locally conformally K\"ahler}, or LCK, if
\[
 \mathrm{d}\omega=\lambda\wedge\omega,
 \qquad \mathrm{d}\lambda=0
\]
for a real one-form $\lambda$, called its \emph{Lee form}. It is \emph{strictly LCK} if $\lambda$ is not exact. A strictly LCK metric is \emph{Vaisman} if its Lee form is parallel for its Levi--Civita connection.
\end{definition}

The strictness convention excludes the K\"ahler case from the word Vaisman. On a compact connected manifold a nonzero parallel one-form cannot be exact: the differential of a global potential vanishes at a maximum, and a parallel form has constant length. Thus a nonzero parallel Lee form is automatically nowhere zero and nonexact.

We use $\lambda$ for the Lee form so that $\theta$ continues to denote a $(1,1)$-background. The map $a\mapsto a\wedge\omega$ is injective on one-forms for $n\geq2$: wedging with $\omega^{n-2}$ reduces this assertion to the symplectic isomorphism $a\mapsto a\wedge\omega^{n-1}$. Thus the Lee form is determined by $\omega$.

\Needspace{5\baselineskip}
Locally write $\lambda=\mathrm{d}f$. Then
\[
 \Omega=\mathrm{e}^{-f}\omega>0,\qquad \mathrm{d}\Omega=0.
\]
The same construction on a covering $\pi:\widetilde X\to X$ on which $\pi^*\lambda=\mathrm{d}F$ gives a K\"ahler form
\begin{equation}\label{eq:lck-cover}
 \widetilde\omega=\mathrm{e}^{-F}\pi^*\omega.
\end{equation}
For a deck transformation $g$, the difference $g^*F-F=c_g$ is constant, and
\[
 g^*\widetilde\omega=\chi(g)\widetilde\omega,
 \qquad \chi(g)=\mathrm{e}^{-c_g}.
\]
This positive character describes the change of scale. Forms transforming by $\chi$ may equivalently be viewed as forms with values in the flat weight line bundle. This is the sense in which an LCK form has a closed interpretation on the cover \cite[Section 3.4]{OV24LCK}.

The conformal change $\omega'=\mathrm{e}^a\omega$ changes the Lee form to $\lambda+\mathrm{d}a$. In particular, the de Rham class $[\lambda]$ depends only on the conformal class. If $\lambda$ is exact, the construction above gives a K\"ahler form on $X$ itself.

\section{Gauduchon representatives and the Lee class}

Set
\[
 \lambda^{\mathrm c}
 =\frac{\mathrm i}{4\pi}(\lambda^{0,1}-\lambda^{1,0}).
\]
Thus $\lambda^{\mathrm c}=\dc f$ wherever $\lambda=\mathrm{d}f$. The closed real $(1,1)$-form
\begin{equation}\label{eq:lck-lee-curvature}
 \kappa_\lambda:=\dc\lambda=-\mathrm{d}\lambda^{\mathrm c}
\end{equation}
will connect the Lee class with ordinary Bott--Chern cohomology. It is de Rham exact, but need not be $\ddc$-exact globally.

\begin{proposition}\label{prop:lck-gauduchon}
An LCK metric is Gauduchon if and only if its Lee form is co-closed for that metric. Its conformal Gauduchon representative therefore has harmonic Lee form. Moreover, for any LCK metric,
\begin{equation}\label{eq:lck-lee-degree}
 \int_X\kappa_\lambda\wedge\omega^{n-1}
 =(n-1)\int_X\lambda\wedge\lambda^{\mathrm c}\wedge\omega^{n-1}.
\end{equation}
If $\lambda\not\equiv0$, the right side is strictly positive.
\end{proposition}
\begin{proof}
On a local conformal K\"ahler chart,
\[
 \dc(\omega^{n-1})
 =(n-1)\lambda^{\mathrm c}\wedge\omega^{n-1}.
\]
The pointwise Hermitian Hodge-star identity identifies
$\lambda^{\mathrm c}\wedge\omega^{n-1}$ with a fixed nonzero constant times $*\lambda$. Consequently $\ddc(\omega^{n-1})=0$ is equivalent to $\mathrm{d}*\lambda=0$, or $\mathrm{d}^*\lambda=0$. Since $\mathrm{d}\lambda=0$, co-closedness is equivalent to harmonicity. Existence of the conformal representative follows from \cref{thm:gauduchon}; conformal changes preserve the LCK condition.

Stokes' theorem and $\mathrm{d}(\omega^{n-1})=(n-1)\lambda\wedge\omega^{n-1}$ give
\[
 \begin{split}
 -\int_X\mathrm{d}\lambda^{\mathrm c}\wedge\omega^{n-1}
 &=-\int_X\lambda^{\mathrm c}\wedge\mathrm{d}(\omega^{n-1})\\
 &=(n-1)\int_X\lambda\wedge\lambda^{\mathrm c}\wedge\omega^{n-1}.
 \end{split}
\]
For a real one-form,
$\lambda\wedge\lambda^{\mathrm c}
=\mathrm i\lambda^{1,0}\wedge\lambda^{0,1}/(2\pi)$ is semipositive. Its trace against $\omega$ is positive wherever $\lambda\ne0$, proving the last assertion.
\end{proof}

There is a linear map
\begin{equation}\label{eq:lck-lee-bc-map}
 H^1_{\mathrm{dR}}(X,\R)\longrightarrow H^{1,1}_{\mathrm{BC}}(X,\R),
 \qquad [\lambda]\longmapsto\{\kappa_\lambda\}.
\end{equation}
Indeed, replacing $\lambda$ by $\lambda+\mathrm{d}a$ replaces $\kappa_\lambda$ by $\kappa_\lambda-\ddc a$. Pairing this class with any Gauduchon form $\gamma^{n-1}$ is therefore well defined. For a strictly LCK structure, its own conformal Gauduchon representative gives a strictly positive pairing by \eqref{eq:lck-lee-degree}. In particular $\{\kappa_\lambda\}\ne0$, although its de Rham image is zero.

\begin{corollary}\label{cor:lck-no-closed-big-class}
A compact strictly LCK manifold admits no big ordinary Bott--Chern class. In particular it lies outside Fujiki's class $\mathcal C$.
\end{corollary}
\begin{proof}
If there were a big closed class, the modification criterion recalled in \cref{chap:bott-chern-loci} would give a K\"ahler modification $\pi:Y\to X$. The form $\pi^*\kappa_\lambda$ is closed and de Rham exact. The K\"ahler $\partial\bar\partial$-lemma gives a smooth real function $h$ on $Y$ with $\pi^*\kappa_\lambda=\ddc h$. Push forward this equality as currents. Since $\pi$ has degree one, $\pi_*\pi^*\kappa_\lambda=\kappa_\lambda$, and pushforward commutes with $\ddc$. Thus $\kappa_\lambda=\ddc(\pi_*h)$, so its Bott--Chern class vanishes. One may use currents to compute this class; alternatively, taking the trace of this equation against a smooth metric and using elliptic regularity shows that the distribution $\pi_*h$ is smooth. This contradicts the positive Gauduchon pairing above. The same modification criterion proves the last assertion.
\end{proof}

For the Hopf example below, $\kappa_\lambda=\alpha$ is smooth semipositive and therefore nef. It has empty non-nef locus, whereas its non-K\"ahler locus is all of $X$, by the non-big convention in \cref{chap:bott-chern-loci}. These closed-class statements coexist with the strictly positive nonclosed background $\omega$ used for its Monge--Amp\`ere equations.

\begin{proposition}\label{prop:lck-balanced-obstruction}
If $X$ has a strictly LCK structure and $\{\kappa_\lambda\}$ is pseudoeffective, then $X$ admits no balanced metric.
\end{proposition}
\begin{proof}
Let $T=\kappa_\lambda+\ddc u\geq0$ represent the class. It is a nonzero closed current, since the class is nonzero. It is also de Rham exact. If $b$ were balanced, then $\mathrm{d}(b^{n-1})=0$, so distributional Stokes would give
\[
 \int_XT\wedge b^{n-1}=0.
\]
The mass of a nonzero positive $(1,1)$-current against a smooth strictly positive $(n-1,n-1)$-form is strictly positive, a contradiction.
\end{proof}

Thus the following conjecture of Ornea--Verbitsky is directly accessible to the Gauduchon duality framework \cite{OV25Balanced}.

\begin{conjecture}\label{conj:lck-lee-psef}
For every compact strictly LCK manifold, the ordinary Bott--Chern class $\{\kappa_\lambda\}$ is pseudoeffective.
\end{conjecture}

The positivity in \eqref{eq:lck-lee-degree} tests the class against the LCK metric in its Gauduchon gauge. Pseudoeffectivity requires nonnegative pairing with \emph{every} Gauduchon metric, by \cref{thm:lamari}. This is the precise additional statement needed for this approach. It is independent of any upper bound for nonlinear Monge--Amp\`ere masses.

\section{Ordinary and twisted potentials}

For the fixed closed one-form $\lambda$, define
\[
 \mathrm{d}_\lambda=\mathrm{d}-\lambda\wedge,
 \qquad \dc_\lambda=\dc-\lambda^{\mathrm c}\wedge,
 \qquad D_\lambda=\mathrm{d}_\lambda\dc_\lambda.
\]
On a chart where $\lambda=\mathrm{d}f$, these operators satisfy
\begin{equation}\label{eq:lck-gauge-operators}
 \mathrm{d}_\lambda=\mathrm{e}^f\mathrm{d}\mathrm{e}^{-f},
 \qquad \dc_\lambda=\mathrm{e}^f\dc \mathrm{e}^{-f},
 \qquad D_\lambda u=\mathrm{e}^f\ddc(\mathrm{e}^{-f}u).
\end{equation}
The identities also show that $D_\lambda u$ is a real $(1,1)$-form for real $u$, and that $\mathrm{d}_\lambda D_\lambda u=0$.

\begin{proposition}\label{prop:lck-two-potentials}
Let $u$ be smooth and real. Then
\[
 \mathrm{d}(\omega+\ddc u)-\lambda\wedge(\omega+\ddc u)
 =-\lambda\wedge\ddc u,
\]
whereas
\[
 \mathrm{d}(\omega+D_\lambda u)
 =\lambda\wedge(\omega+D_\lambda u).
\]
Consequently, an ordinary potential preserves the LCK equation with the same Lee form precisely when $\lambda\wedge\ddc u=0$. A twisted potential always preserves that equation, and gives an LCK metric whenever $\omega+D_\lambda u>0$.
\end{proposition}
\begin{proof}
The first identity follows from $\mathrm{d}\ddc u=0$ and $\mathrm{d}\omega=\lambda\wedge\omega$. The second follows from $\mathrm{d}_\lambda\omega=0$ and \eqref{eq:lck-gauge-operators}.
\end{proof}

On the K\"ahler covering, a twisted potential becomes the automorphic function $\widetilde u=\mathrm{e}^{-F}\pi^*u$, and
\begin{equation}\label{eq:lck-lift-potential}
 \mathrm{e}^{-F}\pi^*(\omega+D_\lambda u)
 =\widetilde\omega+\ddc\widetilde u.
\end{equation}
The numerator in the corresponding Monge--Amp\`ere measure transforms by $\chi^n$. Its ratio to $\widetilde\omega^n$ is invariant and descends to $X$. A scalar measure on $X$ is obtained after fixing the metric gauge, rather than by integrating the automorphic top form over the noncompact cover.

\Needspace{5\baselineskip}
The distinction also affects normalization. Expanding the twisted Hessian gives
\begin{equation}\label{eq:lck-twisted-expansion}
 D_\lambda u
 =\ddc u-\mathrm{d}u\wedge\lambda^{\mathrm c}
       -\lambda\wedge\dc u
       +u\bigl(\kappa_\lambda+\lambda\wedge\lambda^{\mathrm c}\bigr).
\end{equation}
In particular,
$D_\lambda1=\kappa_\lambda+\lambda\wedge\lambda^{\mathrm c}$ need not vanish. Adding a constant changes the metric. The ordinary convention $\sup_Xu=0$ therefore cannot be imposed as a harmless normalization for a twisted equation.

In a chosen gauge the real twisted Bott--Chern group in bidegree $(1,1)$ is
\[
 H^{1,1}_{\mathrm{BC},\lambda}(X,\R)
 =\frac{\{\eta\in C^\infty(X,\Lambda^{1,1}_{\R}):
                       \mathrm{d}_\lambda\eta=0\}}
        {D_\lambda C^\infty(X,\R)}.
\]
This is cohomology with values in the weight line bundle \cite[Section 18.5]{OV24LCK}. It is different both from ordinary Bott--Chern cohomology and from the quotient of all backgrounds by ordinary $\ddc$-terms in \cref{chap:nonclosed-bclasses}.

\section{The Vaisman Hopf example}

\begin{proposition}\label{prop:lck-hopf-vaisman}
The Gauduchon form $\omega=\alpha+\beta$ in \cref{chap:hopf} is Vaisman, with Lee form $\lambda=-\mathrm{d}s$. It satisfies
\[
 \kappa_\lambda=\alpha,
 \qquad D_\lambda1=\omega.
\]
In particular, the bounded mass property fails for an ordinary potential theory whose background is simultaneously Gauduchon and Vaisman.
\end{proposition}
\begin{proof}
On the covering $s=\log|z|^2$, $\alpha=\ddc s$ and $\beta=\mathrm{d}s\wedge\dc s$. Direct differentiation gives
\[
 \mathrm{d}\omega=-\mathrm{d}s\wedge\alpha
                   =-\mathrm{d}s\wedge\omega,
 \qquad \mathrm{e}^s\omega=\ddc(\mathrm{e}^s)=\ddc|z|^2.
\]
Although $s$ is not a function on the quotient, $\mathrm{d}s$ descends. Its integral around the loop induced by $z\mapsto \mathrm{e}^{-1}z$ is $-2$, so the Lee form is not exact. Up to a constant factor, the associated Riemannian metric on the cover is
\[
 \frac{|\mathrm{d}z|^2}{|z|^2}
 =\mathrm{d}t^2+g_{S^5},\qquad t=\log|z|.
\]
The one-form $\mathrm{d}t$ is parallel for this product metric, proving the Vaisman condition. Finally,
$\kappa_\lambda=\ddc s=\alpha$ and
$\lambda\wedge\lambda^{\mathrm c}=\beta$.
The unbounded masses follow from \cref{thm:hopf-unbounded}.
\end{proof}

The last statement concerns the full family $\omega+\ddc u$. It makes no assertion that the particular metrics producing unbounded mass remain Vaisman or have the same Lee form. This distinction is essential when comparing it with volume theorems formulated within the Vaisman class.

\section{Prescribed volume and the Gauduchon condition}

The natural geometric uniqueness question fixes a Lee class and a twisted Bott--Chern class. If the Lee forms themselves differ, one must first specify the identification of the twisted groups. Multiplication by $\mathrm{e}^{-a}$ intertwines the complexes for $\lambda+\mathrm{d}a$ and $\lambda$:
\[
 \mathrm{d}_\lambda(\mathrm{e}^{-a}\eta)
 =\mathrm{e}^{-a}\mathrm{d}_{\lambda+\mathrm{d}a}\eta,
 \qquad
 D_\lambda(\mathrm{e}^{-a}u)=\mathrm{e}^{-a}D_{\lambda+\mathrm{d}a}u.
\]
We fix this identification when comparing classes. The additive constant in $a$ corresponds to a constant rescaling of the weight-bundle identification.

\begin{proposition}\label{prop:lck-coupled-equation}
Let $\omega_1$ be an LCK Gauduchon metric with Lee form $\lambda$. Suppose $\omega_2$ is an LCK metric with Lee form $\lambda+\mathrm{d}a$ and that, under the specified identification,
\[
 \mathrm{e}^{-a}\omega_2=\omega_1+D_\lambda u=:A>0.
\]
Then $\omega_2$ is Gauduchon and $\omega_2^n=\omega_1^n$ if and only if
\begin{equation}\label{eq:lck-coupled-equation}
 \begin{cases}
  A^n=\mathrm{e}^{-na}\omega_1^n,\\
  \ddc\bigl(\mathrm{e}^{(n-1)a}A^{n-1}\bigr)=0.
 \end{cases}
\end{equation}
\end{proposition}
\begin{proof}
Substitute $\omega_2=\mathrm{e}^aA$ into the equal-volume equation and into $\ddc(\omega_2^{n-1})=0$.
\end{proof}

Thus this problem couples a twisted Monge--Amp\`ere equation to a conformal Gauduchon equation. The scalar Hermitian equation of \cref{chap:smooth-equations} does not impose the second equation and uses $\ddc u$ in place of $D_\lambda u$. The prescribed-volume Gauduchon theorem, \cref{thm:stw-volume}, also does not assert that its solution has an LCK structure or the specified twisted class.

\begin{problem}\label{q:lck-gauduchon-volume}
Does a compact LCK manifold have at most one Gauduchon LCK metric with a specified Lee class, specified twisted Bott--Chern class and specified smooth positive volume form?
\end{problem}

This is the uniqueness problem asked in \cite[Question 47.56(3)]{OV24LCK}, with the class identification made explicit. The question concerns metrics: uniqueness of a scalar representative $u$ would additionally require accounting for the kernel of $D_\lambda$. The degenerate-density counterexamples in \cref{chap:degenerate-uniqueness} concern ordinary potentials and do not decide this question.

\subsection{Invariant volumes on Vaisman manifolds}
\label{subsec:lck-vaisman-volume}

The Lee and anti-Lee flows preserve a Vaisman metric. For an invariant prescribed volume, uniqueness in the smooth Hermitian equation forces its potential to be constant along these flows. Transverse Stokes integration then fixes the normalization scalar. Changing the Lee form by the differential of the potential recovers a Vaisman metric with the prescribed volume.

\Needspace{10\baselineskip}
\begin{theorem}\label{thm:lck-vaisman-volume}\index{Vaisman metric!prescribed volume}
Let $(X,\omega,\lambda)$ be a compact Vaisman manifold, and let $\mu$ be a smooth positive volume form invariant under both the Lee field $\lambda^\sharp$ and the anti-Lee field $J\lambda^\sharp$. Assume
\[
 \int_X\mu=\int_X\omega^n.
\]
There is a unique Vaisman metric $\omega'$ with Lee class $[\lambda]$ and $(\omega')^n=\mu$.
\end{theorem}

We first describe the transverse forms and their integration. Let $g$ be the metric of $\omega$. The Lee form has constant squared norm $s=|\lambda|_g^2>0$. Write
\[
 A=s^{-1}\lambda^\sharp,\qquad B=JA,\qquad
 \tau=\frac{s}{4\pi},\qquad \eta=\kappa_\lambda.
\]
Thus $\lambda(A)=1$ and $\lambda(B)=0$. Since $\nabla\lambda=0$, the identity \eqref{eq:lck-complex-connection} gives
\begin{equation}\label{eq:vaisman-normalized-form}
 \begin{split}
 \tau\omega&=\eta+\lambda\wedge\lambda^{\mathrm c},\\
 \eta(U,JV)&=\frac1{4\pi}\bigl(sg(U,V)-\lambda(U)\lambda(V)-\lambda(JU)\lambda(JV)\bigr).
 \end{split}
\end{equation}
In particular, $\eta$ is a closed semipositive $(1,1)$-form with kernel $\mathcal F=\langle A,B\rangle$. To obtain the displayed identity, use $\lambda^{\mathrm c}=-\lambda\circ J/(4\pi)$ and antisymmetrize $\nabla\lambda^{\mathrm c}$; the terms containing $\nabla\lambda$ vanish.

The field $A$ is parallel, and the same connection formula yields
\[
 \nabla_U B=\tfrac12\bigl(-JU+\lambda(U)B+\lambda(JU)A\bigr),
 \qquad \nabla_AJ=0.
\]
The first identity makes $B$ Killing and gives $[A,B]=0$. The second makes $A$ holomorphic; integrability of $J$ then makes $B=JA$ holomorphic. Their flows preserve $g$, $J$, $\lambda$ and $\eta$. Consequently $\mathcal F$ is a holomorphic foliation and $\eta$ descends to a K\"ahler form on each local transverse space.

A smooth function $v$ is \emph{basic} for $\mathcal F$ if $Av=Bv=0$. For such a function, $\mathrm{d}v$, $\dc v$ and $\ddc v$ are horizontal and invariant under both flows; Cartan's formula gives the assertion for $\ddc v$. Put $m=n-1$ and $\eta_v=\eta+\ddc v$. Then
\begin{equation}\label{eq:vaisman-basic-mass}
 \int_X\eta_v^m\wedge\lambda\wedge\lambda^{\mathrm c}
 =\int_X\eta^m\wedge\lambda\wedge\lambda^{\mathrm c}.
\end{equation}
Indeed, $\eta_v^m-\eta^m=\ddc v\wedge T$, where $T=\sum_{j=0}^{m-1}\eta_v^j\wedge\eta^{m-1-j}$ is closed and horizontal. Since $\mathrm{d}(\lambda\wedge\lambda^{\mathrm c})=\lambda\wedge\eta$, the extra term in $\mathrm{d}(\dc v\wedge T\wedge\lambda\wedge\lambda^{\mathrm c})$ vanishes: it has transverse degree $2m+1$. Stokes' theorem proves \eqref{eq:vaisman-basic-mass}.

\Needspace{12\baselineskip}
The following transverse maximum principle will identify the metric after its scale has been fixed.

\begin{proposition}\label{prop:lck-basic-uniqueness}
Let $\mathcal F$ be a foliation on a connected compact manifold $M$ with transverse complex dimension $m\geq1$. Let $\eta_0,\eta_1$ be transverse K\"ahler forms and let $v$ be a smooth real basic function. Compute $\ddc v$ in transverse complex coordinates. Suppose
\[
 \eta_1=\eta_0+\ddc v,
 \qquad \eta_1^m=\eta_0^m.
\]
Then $v$ is constant and $\eta_1=\eta_0$.
\end{proposition}
\begin{proof}
The function $v$ is constant on each local plaque. On every transverse coordinate chart it satisfies
\[
 \ddc v\wedge\sum_{j=0}^{m-1}\eta_1^j\wedge\eta_0^{m-1-j}=0.
\]
This is a linear elliptic equation with no zeroth-order term. Let $c=\max_Mv$. At a point where $v=c$, the strong maximum principle on a connected transverse chart makes $v$ equal to $c$ throughout that chart. Basicness makes it equal to $c$ on the corresponding foliation chart. Hence $\{v=c\}$ is nonempty, open and closed. Connectedness proves the assertion.
\end{proof}

\begin{proof}[of \cref{thm:lck-vaisman-volume}]
Multiplying $\omega$ by $\tau$ and $\mu$ by $\tau^n$ preserves the assertion, so we may assume $\tau=1$. Thus $\omega=\eta+\lambda\wedge\lambda^{\mathrm c}$. Write $\mu=f\omega^n$; the positive smooth function $f$ is basic.

\emph{Existence.} By \cref{thm:tw-smooth}, there are a unique normalized smooth function $u$ and a scalar $b$ such that
\[
 (\omega+\ddc u)^n=\mathrm{e}^b f\omega^n,
 \qquad \omega+\ddc u>0,\qquad \sup_Xu=0.
\]
Pullback by either distinguished flow preserves the equation and the normalization. Uniqueness makes $u$ basic. Hence $\eta_u=\eta+\ddc u$ is a transverse K\"ahler form, and
\[
 (\omega+\ddc u)^n=n\eta_u^m\wedge\lambda\wedge\lambda^{\mathrm c}.
\]
Integrating and using \eqref{eq:vaisman-basic-mass} and $\int_X\mu=\int_X\omega^n$ gives $b=0$.

Set $\lambda_u=\lambda-\mathrm{d}u$ and
\begin{equation}\label{eq:vaisman-reconstructed-metric}
 \omega_u'=\eta_u+\lambda_u\wedge\lambda_u^{\mathrm c}.
\end{equation}
Here $\eta_u=\kappa_{\lambda_u}$ has kernel $\mathcal F$, while $\lambda_u(A)=1$ and $\lambda_u(B)=0$. These facts give $\omega_u'>0$. Direct differentiation gives $\mathrm{d}\omega_u'=\lambda_u\wedge\omega_u'$. Basicness of $u$ makes $A$ a holomorphic Killing field for the new metric $g_u'$, and \eqref{eq:vaisman-reconstructed-metric} gives
\[
 g_u'(A,\cdot)=\frac1{4\pi}\lambda_u.
\]
The closedness of $\lambda_u$ makes $\nabla^{g_u'}A$ symmetric, whereas the Killing equation makes it skew-symmetric. Thus $A$ and $\lambda_u$ are parallel, so $\omega_u'$ is Vaisman with Lee class $[\lambda]$.

Finally, $\eta_u^m$ has top transverse degree. Every term in $\lambda_u\wedge\lambda_u^{\mathrm c}-\lambda\wedge\lambda^{\mathrm c}$ contains a horizontal one-form. Therefore
\[
 (\omega_u')^n
 =n\eta_u^m\wedge\lambda\wedge\lambda^{\mathrm c}
 =(\omega+\ddc u)^n=\mu.
\]

\emph{Uniqueness.} Let $\widehat\omega$ be another Vaisman solution, with Lee form $\widehat\lambda=\lambda+\mathrm{d}v$. Use hats for its $s,\tau,A,B,\eta$. The form $\eta+\widehat\eta$ is exact and semipositive. Its top power has zero integral and is nonnegative, so it vanishes everywhere. The kernel of a sum of two semipositive Hermitian forms is the intersection of their kernels. Since both kernels are complex lines, they coincide. Thus both metrics have the same foliation $\mathcal F$.

The nowhere vanishing holomorphic field $A-\mathrm iB$ trivializes the holomorphic tangent line of $\mathcal F$. The analogous field $\widehat A-\mathrm i\widehat B$ is a holomorphic multiple of it; compactness makes that multiple constant. Write $\widehat A=aA+bB$, with $a,b\in\mathbb R$. Since $\widehat\lambda(\widehat A)=1$ and $\widehat\lambda(\widehat B)=0$,
\[
 \widehat A v=1-a,\qquad \widehat B v=b.
\]
Both fields are complete and $v$ is bounded. These derivatives can be constant only if $a=1$ and $b=0$. Hence $\widehat A=A$, $\widehat B=B$, and $v$ is basic.

\Needspace{7\baselineskip}
The contraction formula following from \eqref{eq:vaisman-normalized-form} is
\begin{equation}\label{eq:vaisman-volume-contraction}
 \iota_B\iota_A\widehat\omega^n
 =\frac{n}{4\pi\widehat\tau^n}\widehat\eta^m,
 \qquad
 \iota_B\iota_A(\omega_u')^n=\frac{n}{4\pi}\eta_u^m.
\end{equation}
Equality of their volume forms gives $\widehat\eta^m=\widehat\tau^n\eta_u^m$. Since $\widehat\eta=\eta-\ddc v$ and $v$ is basic, integrating against $\lambda\wedge\lambda^{\mathrm c}$ and applying \eqref{eq:vaisman-basic-mass} twice gives $\widehat\tau^n=1$. Thus $\widehat\tau=1$ and $\widehat\eta^m=\eta_u^m$. Their difference is $-\ddc(v+u)$, so \cref{prop:lck-basic-uniqueness} gives $v+u$ constant. It follows that $\widehat\lambda=\lambda_u$, $\widehat\eta=\eta_u$ and $\widehat\omega=\omega_u'$.
\end{proof}

\Needspace{19\baselineskip}
\section{Chern--Ricci flatness and Vaisman rigidity}

Return to an LCK metric $\omega$ with Lee form $\lambda$. In Gauduchon gauge, $\lambda$ is harmonic. Chern--Ricci flatness then expresses the Riemannian Ricci tensor in terms of that form and its covariant derivative. Combining the contracted Bianchi identity with the Bochner formula forces the derivative to vanish.

Use $\operatorname{Ric}^{\mathrm{Ch}}=\operatorname{Ric}/(2\pi)$ for the normalized Chern--Ricci form representing $c_1^{\mathrm{BC}}(X)$, with $\operatorname{Ric}$ as in \cref{subsec:chern-ricci-maximal-time}.

\begin{theorem}\label{thm:lck-chern-ricci-rigidity}\index{Chern--Ricci curvature!LCK rigidity}
A compact strictly LCK metric that is both Gauduchon and Chern--Ricci flat is Vaisman.
\end{theorem}

\begin{proof}
We express the Ricci tensor of the Riemannian metric in terms of the Lee form, its covariant derivative and its divergence, and then compare two integral identities for this derivative, one from the contracted Bianchi identity and one from the Bochner formula. Norms, traces and contractions are taken with respect to $g$. For real tangent vectors $U,V$, put
\[
 A=\lambda^\sharp,\qquad B=JA,\qquad
 S=\nabla\lambda,\qquad
 S^+(U,V)=\tfrac12\bigl(S(U,V)+S(JU,JV)\bigr).
\]
The tensor $S$ is symmetric because $\mathrm{d}\lambda=0$, and $S^+$ is its orthogonal projection onto the $J$-invariant symmetric tensors. Thus $S^-=S-S^+$ and $|S|^2=|S^+|^2+|S^-|^2$. Put $\delta=\mathrm{d}^*\lambda$, so that $\delta=-\operatorname{tr}_gS=-\operatorname{div}A$, and write $J\lambda=-\lambda\circ J$, so that $\lambda^{\mathrm c}=J\lambda/(4\pi)$ and $(J\lambda)^\sharp=B$.

\emph{Step 1: the Ricci tensor.} We first show that
\begin{equation}\label{eq:lck-ricci-real}
 \operatorname{Ric}_g
 =nS^+-(n-1)S+\tfrac12\bigl(\delta-|\lambda|^2\bigr)g
       +\tfrac12\lambda\otimes\lambda
       +\tfrac n2 J\lambda\otimes J\lambda.
\end{equation}
To prove this identity, take a local function $f$ with $\lambda=\mathrm{d}f$ and put $\widetilde g=\mathrm{e}^{-f}g=\mathrm{e}^{2\varphi}g$, where $\varphi=-f/2$. The formula for the Ricci tensor under a conformal change, in real dimension $2n$, is
\[
 \operatorname{Ric}_{\widetilde g}
 =\operatorname{Ric}_g-(2n-2)\bigl(\nabla\mathrm{d}\varphi-\mathrm{d}\varphi\otimes\mathrm{d}\varphi\bigr)
       -\bigl(\operatorname{tr}_g\nabla\mathrm{d}\varphi+(2n-2)|\mathrm{d}\varphi|^2\bigr)g.
\]
Since $\nabla\mathrm{d}\varphi=-S/2$ and $\operatorname{tr}_g\nabla\mathrm{d}\varphi=\delta/2$, this gives
\[
 \operatorname{Ric}_{\widetilde g}
 =\operatorname{Ric}_g+(n-1)S
       +\tfrac{n-1}{2}\lambda\otimes\lambda
       -\tfrac12\bigl(\delta+(n-1)|\lambda|^2\bigr)g.
\]
The local metric $\widetilde g$ is K\"ahler, and its volume form is $\mathrm{e}^{-nf}$ times that of $g$. Since the Chern--Ricci form of $g$ vanishes, the Ricci form of $\widetilde g$, written with the unnormalized operator $\mathrm i\partial\bar\partial$, is
$n\mathrm i\partial\bar\partial f=\tfrac n2\mathrm{d}(J\lambda)$.

Applying \eqref{eq:lck-complex-connection} to
$(\nabla_UJ\lambda)(V)=-S(U,JV)-\lambda((\nabla_UJ)V)$ and antisymmetrizing yields
\[
 \mathrm{d}(J\lambda)(U,JV)
 =2S^+(U,V)-|\lambda|^2g(U,V)
             +\lambda(U)\lambda(V)+J\lambda(U)J\lambda(V).
\]
For a K\"ahler metric, the Ricci tensor is obtained from the Ricci form by evaluation on $(U,JV)$. Comparing the last two displays with the conformal formula proves \eqref{eq:lck-ricci-real}. The calculation is local, but the resulting tensor identity holds globally. Since $\operatorname{tr}_gS^+=\operatorname{tr}_gS=-\delta$, taking traces in \eqref{eq:lck-ricci-real} gives the scalar curvature
\[
 \operatorname{scal}_g=(n-1)\delta-\tfrac{n-1}{2}|\lambda|^2.
\]

\emph{Step 2: the Bianchi identity.} All integrals below are taken with respect to the Riemannian volume of $g$, which we omit from the notation. Since
$S(A,A)=\tfrac12A(|\lambda|^2)$ and $\operatorname{div}A=-\delta$,
\begin{equation}\label{eq:lck-aa-zero}
 \int_X S(A,A)=\frac12\int_X|\lambda|^2\delta.
\end{equation}
By the contracted Bianchi identity, the divergence of the Ricci tensor is half the differential of the scalar curvature. Integrating by parts twice therefore gives
\[
 \int_X\langle\operatorname{Ric}_g,S\rangle
 =-\frac12\int_XA(\operatorname{scal}_g)
 =-\frac12\int_X\operatorname{scal}_g\,\delta.
\]
We substitute \eqref{eq:lck-ricci-real} and the scalar curvature into this identity. Since $S^+$ is the orthogonal projection of $S$ onto the $J$-invariant symmetric tensors and $\langle g,S\rangle=-\delta$, and since \eqref{eq:lck-aa-zero} expresses $\int_X|\lambda|^2\delta$ through $\int_XS(A,A)$, we obtain
\begin{equation}\label{eq:lck-bianchi-integral}
 n\int_X|S^+|^2-(n-1)\int_X|S|^2+\frac{n-2}{2}\int_X\delta^2
 +\frac{4-n}{2}\int_XS(A,A)+\frac n2\int_XS(B,B)=0.
\end{equation}

\emph{Step 3: the Bochner formula.} Pairing the Weitzenb\"ock formula for one-forms with $\lambda$ and integrating gives, since $\mathrm{d}\lambda=0$,
\[
 \int_X\delta^2=\int_X|S|^2+\int_X\operatorname{Ric}_g(A,A).
\]
When \eqref{eq:lck-ricci-real} is evaluated on $(A,A)$, its two quartic terms cancel and the term $J\lambda\otimes J\lambda$ vanishes. Thus
\[
 \operatorname{Ric}_g(A,A)
 =(1-\tfrac n2)S(A,A)+\tfrac n2S(B,B)+\tfrac12|\lambda|^2\delta.
\]
Together with \eqref{eq:lck-aa-zero}, this gives
\begin{equation}\label{eq:lck-bochner-integral}
 \int_X\delta^2=\int_X|S|^2+\frac{4-n}{2}\int_XS(A,A)+\frac n2\int_XS(B,B).
\end{equation}
The last two terms are the same in \eqref{eq:lck-bianchi-integral} and \eqref{eq:lck-bochner-integral}. Eliminating them leaves
\[
 n\int_X|S^+|^2-n\int_X|S|^2+\frac n2\int_X\delta^2=0.
\]
Since $|S|^2-|S^+|^2=|S^-|^2$, this is the asserted identity.
\end{proof}


The hypotheses concern the same metric. Vanishing of $c_1^{\mathrm{BC}}(X)$ gives a Chern--Ricci-flat representative in each Hermitian conformal class: if a smooth real function $h$ satisfies $\operatorname{Ric}^{\mathrm{Ch}}(\omega)=\ddc h$ in this normalization, then
\[
 \operatorname{Ric}^{\mathrm{Ch}}(\mathrm{e}^{h/n}\omega)
 =\operatorname{Ric}^{\mathrm{Ch}}(\omega)-\ddc h=0.
\]
Gauduchon's theorem gives a Gauduchon representative in the same conformal class. The rigidity theorem applies when these representatives coincide. Arranging this coincidence is an approach to the following conjecture.

\begin{conjecture}\label{conj:lck-calabi-yau-vaisman}
A compact complex manifold with $c_1^{\mathrm{BC}}(X)=0$ that admits a strictly LCK metric admits a Vaisman metric.
\end{conjecture}

The equation problem is to arrange Chern--Ricci flatness and the Gauduchon condition within an LCK family. The conjecture asks for one Vaisman metric; it does not assert parallelness of every Lee form on the manifold.

\section{The Lee--Gauduchon cone}
\label{sec:lee-gauduchon-cone}

For the rest of the chapter, let $X$ be an arbitrary compact complex manifold of dimension $n\geq2$. To detect pseudoeffectivity in the image of the Lee map $j_X$ of \eqref{eq:lck-lee-bc-map}, one must pair against every Gauduchon metric. The closed forms $\dc(\gamma^{n-1})$ record these tests in de Rham cohomology. Their natural ambient space is the annihilator of $\ker j_X$, where variation of the metric gives an open convex cone and strict duality detects every nonzero positive class in the image.

\Needspace{12\baselineskip}
Set $V_X=H^1_{\mathrm{dR}}(X,\R)$ and
\[
 K_X=\ker j_X,\qquad
 W_X=K_X^\perp\subset H^{2n-1}_{\mathrm{dR}}(X,\R),
\]
where the annihilator uses the Poincar\'e pairing with a class $a\in H^{2n-1}_{\mathrm{dR}}(X,\R)$
\[
 \langle[\lambda],a\rangle=\int_X\lambda\wedge a.
\]
Finally, let
\[
 \mathcal C_X
 =\{[\lambda]\in V_X:j_X([\lambda])\text{ is pseudoeffective}\}.
\]
The image of $j_X$ can be strictly smaller than the kernel of the Bott--Chern-to-de-Rham map, as \cref{ex:lee-calabi-eckmann} will show.

\begin{lemma}\label{lem:lee-gauduchon-kernel}
The space $K_X$ consists of the de Rham classes of real parts of closed holomorphic one-forms. Moreover,
\[
 \mathcal C_X\cap(-\mathcal C_X)=K_X.
\]
\end{lemma}
\begin{proof}
If $j_X([\lambda])=0$, write $\dc\lambda=\ddc f$ for a smooth real function $f$. Since $\dc\mathrm{d}=-\mathrm{d}\dc$, the form $\lambda+\mathrm{d}f$ is both $\mathrm{d}$-closed and $\dc$-closed. Its $(1,0)$-part is therefore holomorphic and closed: the two equations separate the $\partial$ and $\bar\partial$ derivatives. Thus $\lambda+\mathrm{d}f$ is the real part of the closed holomorphic form $2(\lambda+\mathrm{d}f)^{1,0}$. Conversely, the real part of a closed holomorphic form is closed for both operators, and its image under $j_X$ is zero.

By \cref{prop:psef-cone-compact-base}, the pseudoeffective Bott--Chern cone contains no nonzero line. Thus $j_X(\xi)$ and $-j_X(\xi)$ are both pseudoeffective exactly when $j_X(\xi)=0$. This is the asserted identity $\mathcal C_X\cap(-\mathcal C_X)=K_X$.
\end{proof}

\Needspace{5\baselineskip}
\begin{definition}\label{def:lee-gauduchon-cone}\index{Lee--Gauduchon cone}
The \emph{Lee--Gauduchon cone} is
\[
 \operatorname{LG}(X)
 =\bigl\{[\dc(\gamma^{n-1})]:
                    \gamma\text{ is a Gauduchon metric on }X\bigr\}.
\]
\end{definition}

For a closed one-form $\lambda$ and any smooth real $(n-1,n-1)$-form $\Gamma$ with $\ddc\Gamma=0$, integration by parts gives
\begin{equation}\label{eq:lck-lee-gauduchon-pairing}
 \int_X\lambda\wedge\dc\Gamma
 =\int_X\kappa_\lambda\wedge\Gamma.
\end{equation}
The form $\dc\Gamma$ is closed, and both sides depend only on the relevant cohomology classes. In particular the class on the left annihilates $K_X$.

\begin{proposition}\label{prop:lee-gauduchon-open}
The cone $\operatorname{LG}(X)$ is nonempty, convex and open in $W_X$. Its linear span is $W_X$.
\end{proposition}
\begin{proof}
Let $\mathcal Z_X$ be the real vector space of smooth $(n-1,n-1)$-forms $\Gamma$ satisfying $\ddc\Gamma=0$, and set
\[
 q_X(\Gamma)=[\dc\Gamma].
\]
First we show $q_X(\mathcal Z_X)=W_X$. The inclusion into $W_X$ follows from \eqref{eq:lck-lee-gauduchon-pairing}. If $[\lambda]$ annihilates the image of $q_X$, its pairing with every Gauduchon form is zero. The generalized Lamari criterion, \cref{thm:lamari}, implies that both $\kappa_\lambda$ and $-\kappa_\lambda$ have pseudoeffective Bott--Chern classes. By \cref{lem:lee-gauduchon-kernel}, $[\lambda]\in K_X$. Thus the image and $W_X$ have the same annihilator in $V_X$. The finite-dimensional nondegenerate Poincar\'e pairing proves their equality.

Let $\mathcal Z_X^+$ denote the strictly positive forms in $\mathcal Z_X$. Every such form has a unique positive $(n-1)$st root, by the matrix construction used in the proof of \cref{thm:lamari}. This root is Gauduchon. Hence
\[
 \operatorname{LG}(X)=q_X(\mathcal Z_X^+).
\]
Gauduchon's theorem makes $\mathcal Z_X^+$ nonempty. It is convex and stable under positive scaling.

To see openness directly, choose forms $\Gamma_1,\ldots,\Gamma_r\in\mathcal Z_X$ whose images form a basis of $W_X$. If $\Gamma\in\mathcal Z_X^+$, then $\Gamma+\sum_i t_i\Gamma_i$ remains strictly positive for all sufficiently small real $t_i$. Its image describes a neighborhood of $q_X(\Gamma)$ in $W_X$. This also covers $W_X=\{0\}$, with the empty basis. Since a nonempty open subset spans its ambient vector space, the span assertion follows.
\end{proof}

\Needspace{12\baselineskip}
\begin{theorem}\label{thm:lee-gauduchon-duality}\index{Lee--Gauduchon cone!duality}
The cone $\mathcal C_X$ is closed, and
\[
 \mathcal C_X
 =\{\xi\in V_X:\langle\xi,a\rangle\geq0
                         \text{ for every }a\in\operatorname{LG}(X)\}.
\]
Conversely,
\[
 \operatorname{LG}(X)
 =\{a\in W_X:\langle\xi,a\rangle>0
                         \text{ for every }\xi\in\mathcal C_X\setminus K_X\}.
\]
\end{theorem}
\begin{proof}
The first equality is exactly \cref{thm:lamari} applied to the smooth closed background $\kappa_\lambda$, using \eqref{eq:lck-lee-gauduchon-pairing}. It expresses $\mathcal C_X$ as an intersection of closed half-spaces, so this cone is closed.

If $a=[\dc(\gamma^{n-1})]\in\operatorname{LG}(X)$ and $\xi=[\lambda]\in\mathcal C_X\setminus K_X$, choose $T=\kappa_\lambda+\ddc u\geq0$. This current is nonzero because it represents the nonzero class $j_X(\xi)$. Therefore
\[
 \langle\xi,a\rangle
 =\int_XT\wedge\gamma^{n-1}>0.
\]

For the converse, suppose $a\in W_X\setminus\operatorname{LG}(X)$. Finite-dimensional separation of a point from a nonempty open convex cone gives a nonzero functional $\ell\in W_X^*$ such that
\[
 \ell(a)\leq0,\qquad \ell(b)\geq0
                   \quad(b\in\operatorname{LG}(X)).
\]
Indeed, in the usual separating inequalities the constant can be taken to be zero by scaling elements of the cone. By Poincar\'e duality, $\ell$ is represented by a class $\xi\in V_X$, modulo $K_X$. Its nonzero restriction means $\xi\notin K_X$. The first equality in the theorem implies $\xi\in\mathcal C_X$, contradicting the required strict inequalities for $a$. This proves the second equality. If $\operatorname{LG}(X)=W_X$, there is no point to separate and the same conclusion holds.
\end{proof}

\subsection{Bimeromorphic invariance}

A modification preserves degree-one cohomology and detects pseudoeffectivity of classes pulled back from the base. These two facts identify the cones $\mathcal C_X$; duality then identifies the Lee--Gauduchon cones.

\Needspace{10\baselineskip}
\begin{lemma}\label{lem:lee-gauduchon-modification}
Let $\pi:Y\to X$ be a modification between connected compact complex manifolds, and let $\theta$ be a smooth closed real $(1,1)$-form on $X$. Then
\[
 \{\theta\}\text{ is pseudoeffective}
 \quad\Longleftrightarrow\quad
 \{\pi^*\theta\}\text{ is pseudoeffective}.
\]
The pullback on real Bott--Chern cohomology in bidegree $(1,1)$ is injective.
\end{lemma}
\begin{proof}
By \cref{prop:psef-cone-compact-base}, both pullback and pushforward preserve pseudoeffectivity. The underlying potential pullback is the one used in \cref{lem:modification-lone-pullback}. Since $\pi$ has degree one,
\[
 \pi_*\pi^*\{\theta\}=\{\theta\}.
\]
Indeed, the current identity $\pi_*\pi^*\theta=\theta$ follows by integration on the isomorphism locus, whose complement has smooth volume zero. This left inverse proves injectivity and the reverse implication for pseudoeffectivity.
\end{proof}

The degree-one comparison follows by moving cycles off the exceptional image.

\begin{lemma}\label{lem:lee-gauduchon-blowup-cohomology}
If $\pi:Y\to X$ is a modification between connected compact complex manifolds, then
\[
 \pi^*:H^1_{\mathrm{dR}}(X,\R)\longrightarrow
                    H^1_{\mathrm{dR}}(Y,\R)
\]
is an isomorphism. Its Poincar\'e dual is the isomorphism
$\pi_*:H^{2n-1}_{\mathrm{dR}}(Y,\R)\to H^{2n-1}_{\mathrm{dR}}(X,\R)$.
\end{lemma}
\begin{proof}
For a proper analytic subset $A$ of a connected complex manifold $M$, restriction
$H^1(M,\R)\to H^1(M\setminus A,\R)$ is injective, and is an isomorphism
if $A$ has complex codimension at least two. Indeed, choose a Whitney
stratification of $A$. General position moves every one-cycle off its
strata, which have real codimension at least two. If those codimensions
are at least four, it also moves every bounding two-chain off the strata,
relative to its boundary. Thus the map on first homology is surjective
in the first case and bijective in the second. Dualizing over $\R$ gives
the asserted restriction statements.

Since $X$ is smooth, there is an analytic subset $A\subset X$ of complex
codimension at least two such that $\pi$ is an isomorphism over
$U=X\setminus A$. Restriction from $X$ to $U$ is therefore an isomorphism,
whereas restriction from $Y$ to $\pi^{-1}U$ is injective. Every class on
$\pi^{-1}U\simeq U$ extends to $X$ and hence pulls back to a class on $Y$.
The commutative restriction diagram proves that $\pi^*$ is an isomorphism.
Poincar\'e duality and the projection formula give the assertion for $\pi_*$.
\end{proof}

\Needspace{10\baselineskip}
\begin{theorem}\label{thm:lee-gauduchon-bimeromorphic}
Let $X$ and $Y$ be bimeromorphic connected compact complex manifolds. Under the natural identification of $H^1_{\mathrm{dR}}$ and its Poincar\'e-dual identification of $H^{2n-1}_{\mathrm{dR}}$, their spaces $K_X,W_X$ and cones $\mathcal C_X,\operatorname{LG}(X)$ agree.
\end{theorem}
\begin{proof}
First let $\pi:Y\to X$ be a modification. Holomorphic pullback commutes with $\dc$, so
\[
 j_Y(\pi^*\xi)=\pi^*j_X(\xi).
\]
By \cref{lem:lee-gauduchon-modification}, membership in $\mathcal C_X$ is equivalent to membership of the pullback in $\mathcal C_Y$, and membership in $K_X$ is likewise preserved in both directions. By \cref{lem:lee-gauduchon-blowup-cohomology}, every degree-one class upstairs occurs in this way. Thus the degree-one isomorphism identifies the two cones and kernels. Its dual identifies $W_Y$ with $W_X$, and \cref{thm:lee-gauduchon-duality} gives
\[
 \pi_*\operatorname{LG}(Y)=\operatorname{LG}(X).
\]

For a bimeromorphic map between $X$ and $Y$, resolve its graph to obtain
modifications $p:Z\to X$ and $q:Z\to Y$ from one smooth compact manifold.
The preceding argument applies to both maps and identifies the cones and
kernels via $(q^*)^{-1}p^*$ and its Poincar\'e dual. These identifications
are independent of the resolution: they agree on a common open set where
the original map is an isomorphism, and the injectivity of degree-one
restriction proved above determines each class uniquely. The dual
identifications are consequently canonical as well.
\end{proof}

\subsection{Balanced metrics and the Hopf cone}

The presence of zero in the Lee--Gauduchon cone rules out nonzero pseudoeffective classes in the image of $j_X$. Balanced metrics satisfy this condition. On the Hopf threefold, by contrast, the Lee curvature gives a positive ray and the dual cone is an open half-line.

\begin{corollary}\label{cor:lee-gauduchon-zero}
The following conditions are equivalent:
\[
 0\in\operatorname{LG}(X),\qquad
 \operatorname{LG}(X)=W_X,\qquad
 \mathcal C_X=K_X.
\]
A balanced metric implies these conditions. On a manifold in Fujiki's class $\mathcal C$, one has $W_X=\{0\}$ and $\operatorname{LG}(X)=\{0\}$.
\end{corollary}
\begin{proof}
An open cone in a finite-dimensional vector space that contains zero is the entire space: scale a neighborhood of zero to reach any point. The two dual descriptions in \cref{thm:lee-gauduchon-duality} now prove the equivalences, including the case $W_X=\{0\}$. If $b$ is balanced, its form $b^{n-1}$ is closed and of pure type, hence also $\dc$-closed. It is Gauduchon and represents zero in $\operatorname{LG}(X)$.

On a compact K\"ahler manifold the $\partial\bar\partial$-lemma gives $j_X=0$: every $\kappa_\lambda$ is closed and de Rham exact. Thus $K_X=V_X$ and $W_X=0$. Nonemptiness gives $\operatorname{LG}(X)=\{0\}$. Bimeromorphic invariance extends this conclusion to Fujiki's class $\mathcal C$.
\end{proof}

Consequently, a nonzero pseudoeffective class in the image of $j_X$ obstructs balanced metrics throughout the bimeromorphic class. For LCK manifolds, \cref{conj:lck-lee-psef} would provide such a class.

For the Hopf threefold of \cref{chap:hopf}, the covering description identifies the underlying manifold with $S^1\times S^5$, so $V_X=\R[\lambda]$. The class $j_X([\lambda])=\{\alpha\}$ is nonzero by \eqref{eq:lck-lee-degree}, and is pseudoeffective because $\alpha\geq0$. Negative multiples are not pseudoeffective, since their pairing against the fixed Gauduchon metric is negative. Thus
\[
 K_X=0,\qquad \mathcal C_X=\R_{\geq0}[\lambda],\qquad
 \operatorname{LG}(X)=\{a\in H^5_{\mathrm{dR}}(X,\R):
                                  \langle[\lambda],a\rangle>0\}.
\]
The same manifold has unbounded Monge--Amp\`ere mass by \cref{thm:hopf-unbounded}. The positive ray detected here controls linear cohomological pairings, which do not bound those nonlinear masses.

\subsection{Strongly Gauduchon metrics}
\label{sec:lee-gauduchon-strong}

Vanishing of the Lee--Gauduchon class is weaker than the strongly Gauduchon condition. We give an explicit example, following the distinction in \cite[Remarks 5.6--5.8]{OV25LeeGauduchon}. It also shows why the image of $j_X$ must not be replaced by all de Rham-exact Bott--Chern classes.

\begin{definition}
A Hermitian metric $\gamma$ is \emph{strongly Gauduchon} if $\gamma^{n-1}$ is the $(n-1,n-1)$-component of a smooth closed real $(2n-2)$-form.
\end{definition}

Writing that form as
\[
 \Theta=\gamma^{n-1}+\beta+\overline\beta,
 \qquad \beta\text{ of type }(n,n-2),
\]
shows that the definition is equivalent to $\partial(\gamma^{n-1})$ being $\bar\partial$-exact: its value is $-\bar\partial\beta$. Conversely, this equation and its conjugate give $\mathrm{d}\Theta=0$, because the remaining types exceed dimension $n$. Applying $\bar\partial$ also shows that $\gamma$ is Gauduchon. Every balanced metric is strongly Gauduchon by taking $\Theta=\gamma^{n-1}$.

\begin{proposition}\label{prop:lee-strong-obstruction}
A strongly Gauduchon metric has zero Lee--Gauduchon class. If a connected compact complex manifold admits a nonzero positive de Rham-exact $(1,1)$-current, it has no strongly Gauduchon metric, and hence no balanced metric.
\end{proposition}
\begin{proof}
Put $\Gamma=\gamma^{n-1}$ and use $\beta$ as above. In our normalization,
\[
 \dc\Gamma
 =\frac{\mathrm i}{4\pi}
       (\bar\partial\beta-\partial\overline\beta)
 =\mathrm{d}\!\left(\frac{\mathrm i}{4\pi}
                              (\beta-\overline\beta)\right).
\]
Thus $[\dc\Gamma]=0$.

\Needspace{8\baselineskip}
For the obstruction, suppose $T=\mathrm{d}S\geq0$ is nonzero and $\Theta$ is a closed real form with strictly positive $(n-1,n-1)$-part $\Gamma$. Type considerations and distributional Stokes give
\[
 \int_XT\wedge\Gamma
 =\int_XT\wedge\Theta=0.
\]
The first integral is strictly positive because $T\ne0$ and $\Gamma>0$, a contradiction. The argument uses the full current and a smooth test form; it does not involve a non-pluripolar product.
\end{proof}

\begin{example}\label{ex:lee-calabi-eckmann}
There is a connected compact complex threefold $X$ with
\[
 \operatorname{LG}(X)=\{0\},\qquad V_X=K_X=W_X=0,
\]
but with no strongly Gauduchon metric. Moreover, $X$ has a nonzero pseudoeffective Bott--Chern class whose de Rham image is zero and which is not in the image of $j_X$.
\end{example}
\begin{proof}
Consider the holomorphic action of the additive group $\C$ on
\[
 \widetilde X=(\C^2\setminus\{0\})\times(\C^2\setminus\{0\}),
 \qquad
 t\cdot(z,w)=(e^t z,e^{\mathrm i t}w).
\]
The action is free. Indeed, $e^t=e^{\mathrm i t}=1$ forces $t=0$. It is proper: if both $(z,w)$ and $t\cdot(z,w)$ range in compact subsets of $\widetilde X$, the norms of all four vector components are bounded above and away from zero. Writing $t=a+\mathrm i b$, the ratios of the norms bound $e^a$ and $e^{-b}$ above and away from zero. Thus $a$ and $b$ stay bounded, which proves properness. The free proper holomorphic action has local holomorphic slices and a Hausdorff complex quotient $X$ of dimension three.

Every orbit meets $S^3\times S^3$ exactly once. The required group element is determined by
\[
 a=-\log|z|,\qquad b=\log|w|.
\]
These expressions are smooth, so the slice identifies the quotient smoothly with $S^3\times S^3$. In particular $X$ is compact and connected, and
\[
 H^1_{\mathrm{dR}}(X,\R)=H^2_{\mathrm{dR}}(X,\R)
 =H^5_{\mathrm{dR}}(X,\R)=0.
\]
It follows directly that $V_X=K_X=W_X=0$ and $\operatorname{LG}(X)=\{0\}$: the latter cone is nonempty by Gauduchon's theorem and lies in the zero group $H^5_{\mathrm{dR}}(X,\R)$.

\Needspace{16\baselineskip}
Projectivizing both vector components descends to a holomorphic submersion
\[
 p:X\longrightarrow\mathbf P^1\times\mathbf P^1.
\]
Its differential is surjective because the product projectivization on $\widetilde X$ is a submersion and the group orbits lie in its fibers. Let $\eta$ be the sum of the Fubini--Study forms on the two factors. The smooth form
\[
 \alpha=p^*\eta
\]
is closed, semipositive and of rank two everywhere. It is nonzero and de Rham exact, since $H^2_{\mathrm{dR}}(X,\R)=0$. By \cref{prop:lee-strong-obstruction}, it obstructs a strongly Gauduchon metric.

Its Bott--Chern class is nonzero. Otherwise $\alpha=\ddc f$ for a smooth real function $f$, and integration against any Gauduchon form would give zero, contradicting positivity and $\alpha\ne0$. The class is pseudoeffective because it has the positive representative $\alpha$. Finally $j_X$ has zero domain, so this class is not in its image.
\end{proof}

The example distinguishes two linear obstructions. The Lee--Gauduchon cone detects positive classes obtained from closed one-forms by $j_X$. Strongly Gauduchon metrics are obstructed by every nonzero positive de Rham-exact $(1,1)$-current. The latter family can be larger, even when the first cohomology vanishes. Thus the bimeromorphic invariant above is useful without being a complete test for balanced or strongly Gauduchon geometry.


